\clearpage
\section{Results and Discussion}

\subsection{Methodology}

The benchmark ran on OpenJDK 17, Windows 11 amd64, with 16 visible processors.
Each scenario uses two warmups and five measured runs. Only the simulation loop
is timed, and the median limits the effect of isolated pauses. Seed 42 provides
the same initial state, checked against the sequential baseline.

Workload scaling uses 1000--2500 balls, 600 steps, and 16 workers. Worker
scaling uses 2500 balls, 1000 steps, and 1--17 workers. Speedup is
$S=T_{seq}/T_{parallel}$.

\subsection{Speedup as workload grows}

\begin{table}[htbp]
    \centering
    \small
    \setlength{\tabcolsep}{5pt}
    \begin{tabular}{rlrrr}
        \hline
        Balls & Engine & Seq. [ms] & Parallel [ms] & Speedup \\
        \hline
        1000 & threads  & 151.442 & 112.314 & 1.348 \\
        1000 & executor & 151.442 &  99.559 & 1.521 \\
        1500 & threads  & 216.699 & 129.774 & 1.670 \\
        1500 & executor & 216.699 & 134.122 & 1.616 \\
        2000 & threads  & 346.660 & 160.878 & 2.155 \\
        2000 & executor & 346.660 & 168.247 & 2.060 \\
        2500 & threads  & 450.268 & 192.268 & 2.342 \\
        2500 & executor & 450.268 & 213.525 & 2.109 \\
        \hline
    \end{tabular}
    \caption{Median elapsed time and speedup over 600 steps with 16 workers.}
    \label{tab:workload-speedup}
\end{table}

Both concurrent engines outperform the sequential baseline. Speedup rises with
the workload because useful parallel work increasingly amortizes dispatch and
phase barriers. The executor is fastest at 1000 balls ($1.521\times$); from
1500 balls onward, the threaded engine leads and reaches $2.342\times$ at 2500
balls. Their relative overhead therefore depends on workload granularity.

\begin{figure}[htbp]
    \centering
    \includegraphics[width=0.85\linewidth]{images/speedup-vs-balls.png}
    \caption{Speedup as the number of balls increases.}
    \label{fig:speedup-balls}
\end{figure}

\subsection{Worker scaling and saturation}

The sequential median for this experiment is 747.376 ms. Both concurrent
engines use the same 2500-ball state and are compared against that value.

\begin{table}[htbp]
    \centering
    \small
    \setlength{\tabcolsep}{6pt}
    \renewcommand{\arraystretch}{0.95}
    \begin{tabular}{lrrr}
        \hline
        Engine & Workers & Median [ms] & Speedup \\
        \hline
        threads  & 1  & 1063.492 & 0.703 \\
        threads  & 2  &  667.021 & 1.120 \\
        threads  & 4  &  420.976 & 1.775 \\
        threads  & 8  &  318.388 & 2.347 \\
        threads  & 16 &  333.740 & 2.239 \\
        threads  & 17 &  341.042 & 2.191 \\
        \hline
        executor & 1  & 1063.715 & 0.703 \\
        executor & 2  &  681.587 & 1.097 \\
        executor & 4  &  447.603 & 1.670 \\
        executor & 8  &  352.062 & 2.123 \\
        executor & 16 &  353.325 & 2.115 \\
        executor & 17 &  357.971 & 2.088 \\
        \hline
    \end{tabular}
    \caption{Worker-count scaling over 1000 steps with 2500 balls.}
    \label{tab:worker-speedup}
\end{table}

\begin{figure}[htbp]
    \centering
    \includegraphics[width=0.80\linewidth]{images/speedup-vs-workers.png}
    \caption{Speedup as the worker count changes.}
    \label{fig:speedup-workers}
\end{figure}

With one worker, coordination costs make both concurrent engines slower than
the sequential version. Performance peaks at 8 workers: $2.347\times$ for
threads and $2.123\times$ for the executor. Increasing the count to 16 or 17
does not add useful speedup.

Saturation follows from sequential phases, barriers, memory traffic, and finite
independent work. At 17 workers the machine is also oversubscribed. Concurrency
therefore gives a measurable gain on large workloads, but the best worker count is
workload- and machine-dependent.
